\chapter{Forward Kinematics: the Product of Exponentials}
\label{ch:fk}

Forward kinematics answers: \emph{given joint angles $\theta$, where is the tool?} The URDF
answers it numerically every time \texttt{robot\_state\_publisher} broadcasts TF; this chapter
is what that computation actually is. Reference: MR ch.~4.

\section{The formula}

Choose a home configuration (all $\theta_i = 0$) and record $M \in SE(3)$, the tool frame at
home. For each joint $i$, record its screw axis $\mathcal S_i$ (in the base frame, at home).
Then
\begin{equation}
  T_{\text{base,tool}}(\theta) \;=\;
  e^{[\mathcal S_1]\theta_1}\, e^{[\mathcal S_2]\theta_2} \cdots e^{[\mathcal S_n]\theta_n}\, M .
  \label{eq:poe}
\end{equation}
Each factor ``turns on'' one joint, in order from base to tip, then $M$ places the tool. For a
revolute joint through point $q_i$ with axis $\hat\omega_i$:
$\mathcal S_i = (\hat\omega_i,\; -\hat\omega_i \times q_i)$.

Why PoE and not Denavit--Hartenberg: no frame-assignment ritual, no four-parameter bookkeeping,
axes read straight off the CAD/URDF, and it composes cleanly with everything in
chapter~\ref{ch:rigid}. (DH still appears in older papers and UR's own documentation ---
recognize it, don't adopt it.)

\section{Worked example: the 2R planar arm}
\label{sec:2r}

Two links, lengths $l_1, l_2$, both joints about $\hat z$, arm along $\hat x$ at home.

\textbf{Setup.} $M = \begin{bmatrix} I & (l_1{+}l_2, 0, 0)^\top \\ 0 & 1\end{bmatrix}$;
$\mathcal S_1 = (0,0,1;\; 0,0,0)$ (axis through origin);
joint 2 passes through $q_2 = (l_1, 0, 0)$:
$v_2 = -\hat z \times q_2 = -(0, l_1, 0)$, so $\mathcal S_2 = (0,0,1;\; 0, -l_1, 0)$.

\textbf{Result} (expanding \eqref{eq:poe}, or by trigonometry):
\begin{equation}
  x = l_1\cos\theta_1 + l_2\cos(\theta_1{+}\theta_2), \qquad
  y = l_1\sin\theta_1 + l_2\sin(\theta_1{+}\theta_2) .
  \label{eq:2rfk}
\end{equation}

\textbf{Sanity limits} (always do these): $\theta = (0,0)$ gives $(l_1{+}l_2,\,0)$ — home.
$\theta = (0, 180^\circ)$ gives $(l_1 - l_2,\, 0)$ — folded back. $\theta = (90^\circ, 0)$
gives $(0,\, l_1{+}l_2)$ — straight up.

\begin{example}[Numbers]
$l_1 = l_2 = 1$, $\theta = (30^\circ, 45^\circ)$:
$x = \cos 30^\circ + \cos 75^\circ = 0.8660 + 0.2588 = 1.1248$;
$y = \sin 30^\circ + \sin 75^\circ = 0.5 + 0.9659 = 1.4659$.
\end{example}

\section{The UR5e through this lens}

The UR5e is six revolute joints; its screw axes come from the frame data in
\texttt{ur\_description} (the \texttt{\_kinematics.yaml} + xacro chain you met in
chapter~\ref{ch:standin}). Structurally: joint~1 yaws about the base $\hat z$; joints 2, 3, 4
pitch about parallel $\hat y$-axes (shoulder, elbow, wrist-1 — this parallel-axis triple is
why the arm has an elegant ``elbow-up/elbow-down'' geometry); joint 5 rolls; joint 6 pitches
the flange. When mech's arm arrives, our first analytical act will be writing ITS six screw
axes from the CAD --- the PoE table \emph{is} the kinematic datasheet.

\begin{keyresult}
The URDF is a serialized PoE model: each \texttt{<joint>}'s \texttt{<origin>} + \texttt{<axis>}
encode where $q_i$ and $\hat\omega_i$ sit relative to the parent. FK in
\texttt{robot\_state\_publisher}, planning in MoveIt, and physics in Isaac all consume this
one description --- which is why URDF correctness (chapter 2's guardrails, mech's CAD export)
is load-bearing for everything.
\end{keyresult}

\section{Problems}

\begin{problem}
For the 2R arm with $l_1 = 0.35$, $l_2 = 0.30$ (our interface-spec proportions), compute the
tool position at $\theta = (45^\circ, -30^\circ)$.
\end{problem}
\begin{answer}
$\theta_1{+}\theta_2 = 15^\circ$.
$x = 0.35\cos45^\circ + 0.30\cos15^\circ = 0.2475 + 0.2898 = 0.5373$;
$y = 0.35\sin45^\circ + 0.30\sin15^\circ = 0.2475 + 0.0776 = 0.3251$.
\end{answer}

\begin{problem}
Add a base yaw joint (about $\hat z$ at the origin, but now the 2R plane is vertical: joints
2, 3 about $\hat y$). Write the three screw axes $\mathcal S_1, \mathcal S_2, \mathcal S_3$ at
home (arm along $\hat x$).
\end{problem}
\begin{answer}
$\mathcal S_1 = (0,0,1;\,0,0,0)$;
$\mathcal S_2 = (0,1,0;\, -\hat y \times q_2)$ with $q_2 = 0 \Rightarrow v_2 = 0$ if the
shoulder sits at the origin, else $q_2=(0,0,h)$ gives $v_2 = -\hat y \times (0,0,h) = (-h,0,0)$;
$\mathcal S_3 = (0,1,0;\, -\hat y \times (l_1,0,h)) = (0,1,0;\; -h, 0, l_1)$.
(If you got $v_3 = (\!-h,0,l_1\!)$ check the cross-product order:
$\hat y \times (l_1,0,h) = (h,\,0,\,-l_1)$, negated $\Rightarrow (-h, 0, l_1)$.) \qedhere
\end{answer}

\begin{problem}
Why must the exponentials in \eqref{eq:poe} be multiplied in base-to-tip order? What physical
statement does swapping $e^{[\mathcal S_1]\theta_1}$ and $e^{[\mathcal S_2]\theta_2}$ make?
\end{problem}
\begin{answer}
Each screw is written in the \emph{home} base frame; turning joint 1 moves joint 2's axis
through space, and left-multiplication applies that displacement to everything downstream.
Swapping claims joint 2 rotates about its home axis regardless of joint 1 --- a different
(wrong) machine. Same reason matrix order mattered in Problem~1.1.
\end{answer}

\begin{problem}[bridge to practice]
In RViz you set the UR5e joints to all-zero with \texttt{joint\_state\_publisher\_gui}. The
arm points straight up, not along $\hat x$. What does that tell you about $M$ in
\texttt{ur\_description}, and why does it not matter for anything downstream?
\end{problem}
\begin{answer}
UR's home configuration has the arm vertical --- $M$ encodes that choice. PoE is indifferent:
$M$ is just ``where the tool is at zero.'' Conventions only need to be consistent between the
description and everything that consumes it --- which the URDF guarantees by construction.
\end{answer}
